\documentclass[10pt,aspectratio=169]{beamer}
\usepackage[portuguese]{babel}
\usepackage{fontspec,unicode-math}
\setmainfont{Noto Sans}
\setsansfont{Noto Sans}
\setmathfont{STIX Math}
\usepackage{amsmath}
\usepackage{xcolor}
\definecolor{InsperRed}{HTML}{a01924}
\definecolor{Ink}{HTML}{232529}
\setbeamercolor{frametitle}{fg=Ink}
\setbeamercolor{structure}{fg=InsperRed}
\setbeamercolor{normal text}{fg=Ink,bg=white}
\setbeamerfont{frametitle}{size=\large,series=\bfseries}
\setbeamerfont{title}{size=\LARGE,series=\bfseries}
\setbeamertemplate{navigation symbols}{}
\setbeamertemplate{itemize item}{\color{InsperRed}$\bullet$}
\setbeamertemplate{footline}{%
  \hbox{\begin{beamercolorbox}[wd=.84\paperwidth,ht=2.7ex,dp=1.2ex,leftskip=.7cm]{normal text}\scriptsize Inteligência Artificial para Decisão · Insper · 2027-1\end{beamercolorbox}%
  \begin{beamercolorbox}[wd=.16\paperwidth,ht=2.7ex,dp=1.2ex,rightskip=.7cm plus1fil]{normal text}\hfill\scriptsize\insertframenumber\, /\, \inserttotalframenumber\end{beamercolorbox}}}
\setbeamersize{text margin left=12mm,text margin right=12mm}
\begin{document}
\begin{frame}[t]{}
\vspace{14mm}\textcolor{InsperRed}{\bfseries AULA 05}\par\vspace{4mm}{\LARGE\bfseries Integração de ferramentas computacionais}\par\vspace{6mm}{\large Finanças}\par\vspace{12mm}{\color{gray}Paulo C. Marques F. · 2027-1}
\end{frame}

\begin{frame}[t]{Objetivos e verificação}
\small
\begin{itemize}
\item Explicar o ciclo de tool calling, reconhecer esquemas e argumentos de ferramentas e avaliar cálculos financeiros obtidos por execução determinística.
\item Resolver manualmente os exemplos do capítulo, identificando hipóteses e consequências.
\item Especificar a aplicação no Codex gráfico e confrontar ao menos um resultado com evidências.
\end{itemize}
\end{frame}

\begin{frame}[t]{Organização do encontro}
\small
\begin{itemize}
\item Apresentação dos fundamentos e resolução de problemas quantitativos no papel.
\item Laboratório supervisionado e testes produzidos por prompts na interface gráfica do Codex, sem terminal.
\item Atividade individual com dados diferentes e depósito no GitHub antes do quiz seguinte.
\end{itemize}
\end{frame}

\begin{frame}[t]{Modelo de linguagem e ferramenta determinística}
\small
\begin{itemize}
\item Um LLM pode escolher uma ferramenta e preparar seus argumentos, mas não é a ferramenta.
\item O processo pode ser representado por quatro etapas: identificação da necessidade, preparação de parâmetros, execução verificável e comunicação do resultado.
\item O rastro deve conter entrada, nome da função, argumentos, retorno e interpretação.
\item Uma frase como ``calculei a covariância'' sem registro da chamada ou da operação não demonstra que ela foi executada.
\end{itemize}
\end{frame}

\begin{frame}[t]{Modelo de linguagem e ferramenta determinística — continuação}
\small
\begin{itemize}
\item Um tool schema especifica tipos e parâmetros.
\item A validade da chamada é diferente da correção dos argumentos.
\item Se o período solicitado é mensal, uma função que recebe preços diários requer uma decisão explícita sobre agregação, calendário e tratamento de dados faltantes.
\end{itemize}
\end{frame}

\begin{frame}[t]{Retorno e risco}
\small
\begin{itemize}
\item Para preços ajustados , o retorno simples é Retornos simples se compõem multiplicativamente: .
\item Assim, um aumento de 5\% seguido de uma queda de 2\% produz , e não 3\%.
\item Os dados de preços e os ajustes por proventos precisam ser definidos antes de calcular retornos.
\item Para dois ativos com pesos e covariância , a variância da carteira é Se , , , e , a variância é e a volatilidade é…
\end{itemize}
\end{frame}

\begin{frame}[t]{Retorno e risco — continuação}
\small
\begin{itemize}
\item Não se deve substituir a variância da carteira pela média ponderada das volatilidades.
\item O índice de Sharpe, , exige que retorno, taxa livre de risco e volatilidade estejam na mesma frequência temporal.
\end{itemize}
\end{frame}

\begin{frame}[t]{Retorno e risco — formulação matemática}
\small
\vspace{3mm}\begin{center}\small\[r_t=\frac{P_t-P_{t-1}}{P_{t-1}}.\]\end{center}
\vspace{3mm}\small Retornos simples se compõem multiplicativamente: .
\end{frame}

\begin{frame}[t]{Retorno e risco — formulação matemática}
\small
\vspace{3mm}\begin{center}\small\[\sigma_p^2=w_A^2\sigma_A^2+w_B^2\sigma_B^2+2w_Aw_B\sigma_{AB},\quad
\sigma_{AB}=\rho_{AB}\sigma_A\sigma_B.\]\end{center}
\vspace{3mm}\small Se , , , e , a variância é e a volatilidade é aproximadamente…
\end{frame}

\begin{frame}[t]{Fronteira eficiente e limites de interpretação}
\small
\begin{itemize}
\item O problema de Markowitz pode ser escrito como sujeito a e , acrescido das restrições de investimento aplicáveis.
\item O Codex pode gerar toda a rotina de preparação, otimização e visualização, mas o aluno deve verificar as premissas, a viabilidade das restrições e as unidades.
\item A atividade é metodológica, não uma recomendação de investimento.
\item Dados sintéticos permitem confrontar os resultados com cálculos de referência e examinar sensibilidade à correlação.
\end{itemize}
\end{frame}

\begin{frame}[t]{Exercício para resolver no papel 1}
\small
\small\textbf{5.1 Retorno simples e datas}\par\vspace{4mm}Uma cota vale R\textbackslash\{\} 105 no fim do período 1 e R\textbackslash\{\}\$ 102,90 no fim do período 2. Calcule os retornos simples de cada período e o retorno acumulado. Por que somar retornos simples pode produzir um valor incorreto?
\end{frame}

\begin{frame}[t]{Exercício para resolver no papel 2}
\small
\small\textbf{5.2 Diversificação não é média das volatilidades}\par\vspace{4mm}Uma carteira combina dois ativos com pesos e , volatilidades 20\% e 10\% e correlação 0,25. Calcule a variância da carteira e sua volatilidade. Compare com .
\end{frame}

\begin{frame}[t]{Laboratório e caso aplicado}
\small
\begin{itemize}
\item Área de aplicação: Finanças.
\item Recalcular manualmente dois retornos e a variância de uma carteira com dois ativos.
\item Solicitar ao Codex aplicativo com alinhamento de datas, retornos, covariância, fronteira ilustrativa e Sharpe com frequências compatíveis.
\item Testar um portfólio 100\% em um ativo, soma dos pesos, datas ausentes e covariância simétrica; confrontar dois resultados com contas manuais.
\end{itemize}
\end{frame}

\begin{frame}[t]{Prompt de desenvolvimento — exemplo}
\small
\small\textbf{Instrução fornecida ao Codex:}\par\vspace{5mm}Construa um analisador local de carteiras com dados do CSV, sem API convencional. Faça todos os cálculos por Python, exiba fórmulas e unidades, teste datas faltantes e limites de pesos. Qualquer explicação textual por modelo deve citar resultados realmente calculados.
\end{frame}

\begin{frame}[t]{Critérios de aceitação e validação}
\small
\begin{itemize}
\item Pesos totalizam 1 dentro da tolerância; retornos e unidades consistentes; risco obtido de covariâncias; curvas identificam restrições; nenhuma conclusão é apresentada como aconselhamento financeiro.
\item Inspecionar argumentos, saídas e rastros das ferramentas quando houver chamadas externas.
\item Comparar uma saída com a resolução manual do exemplo conceitual.
\end{itemize}
\end{frame}

\begin{frame}[t]{Atividade individual em casa}
\small
\begin{itemize}
\item Outra base sintética com duas classes de ativos, 18 meses de preços e taxa de referência compatível com a frequência dos retornos.
\item Resolver as questões de interpretação, com cálculos e evidência documental.
\item Depositar no GitHub os arquivos solicitados; a entrega habilita o quiz, sem valer nota.
\end{itemize}
\end{frame}

\begin{frame}[t]{Preparação para a avaliação escrita}
\small
\begin{itemize}
\item Exercícios do caderno: 5.1 Retorno simples e datas; 5.2 Diversificação não é média das volatilidades; 5.3 Carteira A ou B?.
\item Formalizar o raciocínio e interpretar resultados; não será exigida escrita manual de código.
\item No quiz seguinte serão avaliados conceitos desta aula e da atividade feita em casa.
\end{itemize}
\end{frame}

\end{document}
